\begingroup
\renewcommand{\arraystretch}{1.45}
\fancyhead[L]{\small ARITHMETISCH-GEOMETRISCHES MITTEL.}
\Needspace{34\baselineskip}
\AGMArticle{13}
[Die Differentiale erster Ordnung der einzelnen Glieder des vollständigen Algorithmus eines arithmetisch-geometrischen Mittels sind in Artikel 9 auf zweierlei Weise zueinander in Beziehung gesetzt. Die eine Beziehung umfasst nur die Differentiale der Logarithmen der Quotienten jener Größen. Sie ergibt sich aus der wiederholten Anwendung der beiden Relationen $aa=bb+cc$, $\frac{a'}{c'}=\frac{a+b}{a-b}$ und kann unter Verwendung der im vorhergehenden Artikel gewonnenen Ergebnisse zu einer Bestimmung des Differentials des Quotienten zweier zusammengehöriger arithmetisch-geometrischer Mittel erweitert und so dargestellt werden, dass $\frac{1}{cc}\,\mathrm{d}\log\frac{a}{b}$, das mit $\Delta$ bezeichnet werden mag, für jedes positive und negative $n$, wobei ein negativer nachstehender Index $n$ als gleichbedeutend mit einem voranstehenden positiven Index von gleichem absoluten Betrag verstanden wird,
\SourceLeaf{389}
\Needspace{10\baselineskip}
\[
\begin{aligned}&\\&=\frac{1}{2^{n}}\mathbin{.}\frac{1}{a^{\mathrm{n}}a^{\mathrm{n}}}\,\mathrm{d}\log\frac{c^{\mathrm{n}}}{b^{\mathrm{n}}}\\&=\frac{1}{2^{n}}\mathbin{.}\frac{1}{b^{\mathrm{n}}b^{\mathrm{n}}}\,\mathrm{d}\log\frac{c^{\mathrm{n}}}{a^{\mathrm{n}}}\\&=\frac{1}{2^{n}}\mathbin{.}\frac{1}{c^{\mathrm{n}}c^{\mathrm{n}}}\,\mathrm{d}\log\frac{a^{\mathrm{n}}}{b^{\mathrm{n}}}\\&=\frac{\pi}{2}\mathbin{.}\frac{1}{\mathrm{M}(a,b)\mathbin{.}\mathrm{M}(a,c)}\,\mathrm{d}\log\frac{\mathrm{M}(a,c)}{\mathrm{M}(a,b)}\end{aligned}
\]
ist. \Needspace{21\baselineskip}
Die andere Beziehung erstreckt sich auch auf die Differentiale der Logarithmen von Produkten aus zwei Gliedern. Sie folgt aus $b'b'=ab$ und ${}'c{}'c=4ac$ in der Form
\Needspace{10\baselineskip}
\[
\begin{aligned}
\mathrm{d}\log(a^{\mathrm{m}}\mathbin{.} b^{\mathrm{m}})&=\mathrm{d}\log(a^{\mathrm{m}+1}\mathbin{.} b^{\mathrm{m}+1})\\&\quad-\Delta\mathbin{.}2^{m+1}\mathbin{.} c^{\mathrm{m}+1}\mathbin{.} c^{\mathrm{m}+1}\\[3pt]
\mathrm{d}\log({}^{\mathrm{m}}a\mathbin{.}{}^{\mathrm{m}}c)&=\mathrm{d}\log({}^{\mathrm{m}+1}a\mathbin{.}{}^{\mathrm{m}+1}c)\\&\quad+\Delta\mathbin{.}\frac{1}{2^{m+1}}\mathbin{.}{}^{\mathrm{m}+1}b\mathbin{.}{}^{\mathrm{m}+1}b
\end{aligned}
\]
\Needspace{20\baselineskip}
und ergibt, wenn man $m$ bis zur unendlichen Grenze wachsen lässt, die in Artikel 9 aufgestellte Gleichung und die ihr entsprechende, nämlich:
\Needspace{10\baselineskip}
\[
\begin{aligned}
2\,\mathrm{d}\log\mathrm{M}(a,b)&=\mathrm{d}\log(a^{\mathrm{n}}\mathbin{.} b^{\mathrm{n}})\\&\quad+\Delta\{2^{n+1}\mathbin{.} c^{\mathrm{n}+1}\mathbin{.} c^{\mathrm{n}+1}+\ .\ +2^{n+m}\mathbin{.} c^{\mathrm{n}+\mathrm{m}}\mathbin{.} c^{\mathrm{n}+\mathrm{m}}+\ .\}\\[3pt]
2\,\mathrm{d}\log\mathrm{M}(a,c)&=\mathrm{d}\log({}^{\mathrm{n}}a\mathbin{.}{}^{\mathrm{n}}c)\\&\quad-\Delta\{2^{-n-1}\mathbin{.}{}^{\mathrm{n}+1}b\mathbin{.}{}^{\mathrm{n}+1}b+\ .\ +2^{-n-m}\mathbin{.}{}^{\mathrm{n}+\mathrm{m}}b\mathbin{.}{}^{\mathrm{n}+\mathrm{m}}b+\ .\}
\end{aligned}
\]
\Needspace{22\baselineskip}
\par Hieraus kann durch Elimination der Differentiale mithilfe des oben gefundenen Ausdrucks für $\Delta$ die Gleichung
\Needspace{10\baselineskip}
\[
\begin{aligned}
\frac{4}{\pi}\mathrm{M}(a,b)\,\mathrm{M}(a,c)&=.\ .-2^{-\mu+n}\mathbin{.}{}^{\mu-\mathrm{n}}b\mathbin{.}{}^{\mu-\mathrm{n}}b\\&\quad-.\ .-2^{n-1}\mathbin{.} b^{\mathrm{n}-1}\mathbin{.} b^{\mathrm{n}-1}\\
&{}+2^{n}\mathbin{.} a^{\mathrm{n}}\mathbin{.} a^{\mathrm{n}}-2^{n+1}\mathbin{.} c^{\mathrm{n}+1}\mathbin{.} c^{\mathrm{n}+1}\\&\quad-.\ .-2^{n+m}\mathbin{.} c^{\mathrm{n}+\mathrm{m}}\mathbin{.} c^{\mathrm{n}+\mathrm{m}}-.\ .
\end{aligned}
\]
abgeleitet werden, die \textsc{Gauss} neben der im vorigen Artikel wiedergegebenen ersten Gleichung für $\frac{\pi}{2}\frac{\mathrm{M}(a,b)}{\mathrm{M}(a,c)}$ aufgezeichnet hat, ohne den Weg anzudeuten, auf dem sie gefunden worden waren.]
\Needspace{30\baselineskip}
\AGMArticle{14}
\par [Für die totalen Differentiale zweiter Ordnung findet man unmittelbar aus den Ausdrücken für $\Delta$, dass jedes System von Gliedern $a,b,c$ mit gemeinsamem Index die drei Gleichungen
\Needspace{10\baselineskip}
\[
\begin{aligned}
\frac{1}{\Delta}\,\mathrm{d}\log\frac{c}{b}\mathbin{.}\mathrm{d}\log a&=\tfrac12\,\mathrm{d}\left(\frac{1}{\Delta}\,\mathrm{d}\log\frac{c}{b}\right)\\[4pt]
\frac{1}{\Delta}\,\mathrm{d}\log\frac{c}{a}\mathbin{.}\mathrm{d}\log b&=\tfrac12\,\mathrm{d}\left(\frac{1}{\Delta}\,\mathrm{d}\log\frac{c}{a}\right)\\[4pt]
\frac{1}{\Delta}\,\mathrm{d}\log\frac{a}{b}\mathbin{.}\mathrm{d}\log c&=\tfrac12\,\mathrm{d}\left(\frac{1}{\Delta}\,\mathrm{d}\log\frac{a}{b}\right)
\end{aligned}
\]
\Needspace{15\baselineskip}
erfüllt. Multipliziert man darin die Zähler und Nenner unter den zweimal zu differenzierenden Logarithmen der Reihe nach mit $a,b,c$ und zerlegt alle Logarith\SourceLeaf{390}men von Quotienten in Differenzen von Logarithmen, so erkennt man leicht, dass der Ausdruck
\Needspace{10\baselineskip}
\[
\frac{1}{\Delta}\,\mathrm d\log a^{\mathrm n}\mathbin{.}\mathrm d\log b^{\mathrm n}
-\frac12\mathrm d\left(\frac{1}{\Delta}\,\mathrm d\log a^{\mathrm n}b^{\mathrm n}\right)
\]
der mit $D$ bezeichnet werden soll, seinen Wert nicht ändert, wenn man statt
$a^{\mathrm n}$ und $b^{\mathrm n}$ setzt: $a^{\mathrm n}$ und $c^{\mathrm n}$ oder
$b^{\mathrm n}$ und $c^{\mathrm n}$. Nimmt man das Mittel der beiden letzten so entstandenen Ausdrücke und berücksichtigt, dass
\Needspace{10\baselineskip}
\[
\mathrm d\log a^{\mathrm n}+\mathrm d\log b^{\mathrm n}
=2\mathrm d\log b^{\mathrm n+1}
\]
ist, so folgt, dass man in jenem Ausdruck ohne Änderung seines Wertes statt der beiden genannten Größen auch $c^{\mathrm n}$ und $b^{\mathrm n+1}$ setzen kann; ferner, wenn man darin wieder $\mathrm d\log c^{\mathrm n}$ durch
$\frac12\mathrm d\log a^{\mathrm n+1}+\frac12\mathrm d\log c^{\mathrm n+1}$ ersetzt,
dass man mit demselben Ergebnis $a^{\mathrm n+1}$ und $b^{\mathrm n+1}$ statt
$a^{\mathrm n}$ und $b^{\mathrm n}$ setzen kann. \Needspace{23\baselineskip}
Der Wert jenes Differentialausdrucks ist also unabhängig von $n$ und demnach gleich
\Needspace{23\baselineskip}
\[
\begin{aligned}
D&=\displaystyle\frac{1}{\Delta}\,\mathrm d\log b^{\mathrm n}\mathbin{.}\mathrm d\log c^{\mathrm n}
-\frac12\mathrm d\left(\frac{1}{\Delta}\,\mathrm d\log b^{\mathrm n}c^{\mathrm n}\right)\\[6pt]
&=\displaystyle\frac{1}{\Delta}\,\mathrm d\log c^{\mathrm n}\mathbin{.}\mathrm d\log a^{\mathrm n}
-\frac12\mathrm d\left(\frac{1}{\Delta}\,\mathrm d\log c^{\mathrm n}a^{\mathrm n}\right)\\[6pt]
&=\displaystyle\frac{1}{\Delta}\,\mathrm d\log a^{\mathrm n}\mathbin{.}\mathrm d\log b^{\mathrm n}
-\frac12\mathrm d\left(\frac{1}{\Delta}\,\mathrm d\log a^{\mathrm n}b^{\mathrm n}\right)\\[6pt]
&=\displaystyle\mathrm M(a,b)\,\mathrm d\left\{\frac{1}{\Delta}\,\mathrm d\frac{1}{\mathrm M(a,b)}\right\}\\&=\mathrm M(a,c)\,\mathrm d\left\{\frac{1}{\Delta}\,\mathrm d\frac{1}{\mathrm M(a,c)}\right\}
\end{aligned}
\]
\Needspace{22\baselineskip}
\par Setzt man $a$ als unveränderlich voraus, so wird
\Needspace{10\baselineskip}
\[
D=\Delta bbcc=-bb\,\mathrm d\log b=cc\,\mathrm d\log c
\]
und es entsteht die in Artikel 8 aus der Reihenentwicklung abgeleitete Differentialgleichung, die für
$\frac{a}{\mathrm M(a,c)}$ und $\frac{a}{\mathrm M(a,b)}$ als Werte von $\mu$ gilt:
\Needspace{10\baselineskip}
\[
\mathrm{dd}\mu-\frac{\mathrm d\Delta}{\Delta}\mathbin{.}\mathrm d\mu
-\Delta\Delta bbcc\mathbin{.}\mu=0
\]
Aus ihr ergeben sich nach den Untersuchungen in der Abhandlung
`\textit{Determinatio seriei nostrae per aequationem differentialem secundi ordinis}'
auch wieder die Darstellungen durch die Reihen von \textsc{Gauss}:
\Needspace{10\baselineskip}
\[
\begin{gathered}\frac{a}{\mathrm M(a,b)}=F\left(\frac12,\frac12,1,\frac{cc}{aa}\right),\\[7pt]
\frac{a}{\mathrm M(a,c)}=F\left(\frac12,\frac12,1,\frac{bb}{aa}\right)\end{gathered}
\]
\SourceLeaf{391}
und als Spezialfälle der dortigen Gleichungen [90] und [96] der oben in Artikel 12 gefundene Grenzwert von $\mathrm M(1,\varepsilon)\log\frac{4}{\varepsilon}$ und die im vorigen Artikel aufgestellte Differentialgleichung erster Ordnung zwischen
$\mathrm M(a,b)$ und $\mathrm M(a,c)$.]

\clearpage
\endgroup
